\documentclass[10pt,a4paper]{article}
	
	\usepackage{amsmath,amssymb,amsthm,mathtools,amsfonts}
	\usepackage{breqn}
	\usepackage[utf8]{inputenc}
	\usepackage[english]{babel}
	\usepackage[T1]{fontenc}
	\usepackage{graphicx}
	\graphicspath{{./img/}}
	
	%\usepackage{ebgaramond}
	\usepackage{fontspec}
	\setmainfont[Numbers=OldStyle]{EB Garamond}
	\newfontfamily\plantinb[Numbers=OldStyle]{Plantin MT Pro Bold Condensed}
	\newfontfamily\plantin[Numbers=OldStyle]{Plantin MT Pro Semibold}
	
	\usepackage[pdfborder={0 0 0}]{hyperref} %% Ingen firkant rundt om referencer
	\usepackage{multicol}
	\usepackage{tabularx}
	\usepackage{enumitem}% better controls with enumerating
	\usepackage{wrapfig}
		%\setlength{\wrapoverhang}{\marginparwidth}
		%\addtolength{\wrapoverhang}{\marginparsep}
		\setlength{\columnsep}{0pt}
		\setlength{\intextsep}{-10pt}
	\usepackage{caption}
	\usepackage[svgnames]{xcolor}
	\usepackage{dirtytalk}
	\usepackage[modulo]{lineno}
	\usepackage{hyperref}
	\usepackage{tasks}
	\newcommand\svar[1]{\newline\textcolor{red}{\textbf{Svar}\\#1}}
	\newcommand{\qm}[1]{‘‘#1”}
	\newcommand{\sqm}[1]{‘#1’}
	
	\definecolor{hgblue}{RGB}{10, 40, 80}
	\definecolor{hgred}{RGB}{140, 10, 10}
	\definecolor{hggrey}{RGB}{215, 215, 210}
	\definecolor{hggreen}{RGB}{145, 205, 205}
	
	\usepackage{titlesec}
	\titleformat{\subsection}{\color{hgred}\mysectionfont\large}{}{}{}
	\newcommand{\source}[1]{Source: \href{#1}{#1}}
	\usepackage{lipsum}
	\usepackage{comment}
	\usepackage{ulem}
	
	\usepackage{bigfoot}
		\DeclareNewFootnote{a}
			\renewcommand{\thefootnotea}{\fnsymbol{footnotea}}
				\newcommand{\footnotes}[1]{\footnotea[2]{#1}} % here you can specify what symbol you want in footnotes
				\newcommand{\footnoted}[1]{\footnotea[3]{#1}} % for a second footnote on a page
				\MakePerPage{footnotes}
	
	%\reversemarginpar
	\newcommand{\glos}[3]{\dotuline{#1}\marginpar{\scriptsize\textit{#3} #2}}

	\setlength{\parskip}{0.5\baselineskip}
	\setlength{\parindent}{0cm}
	\setlist[enumerate]{label=\alph*),left=20pt}
	%\setlist[enumerate]{leftmargin=\parindent}
	%\setlist[enumerate]{nosep}
	
	%Titling
	\usepackage{titling}
	\setlength{\droptitle}{-12ex}
	\pretitle{\begin{flushleft}\plantinb\huge\color{hgblue}}
	\posttitle{\par\end{flushleft}}
	\preauthor{\begin{flushleft}\Large}
	\postauthor{\end{flushleft}}
	\predate{\begin{flushleft}}
	\postdate{\end{flushleft}}
	
	\setcounter{secnumdepth}{0}
	
	\usepackage{tikz}

\begin{document}
%\pagecolor{FloralWhite}
\title{New Zealand Prime Minister Addresses United Nations General Assembly}
%\date{\vspace{-3em}} % I tilfældet ingen dato
\date{2019}
\author{Jacinda Ardern}
\maketitle
    
\pagestyle{empty}
\thispagestyle{empty}
\linenumbers
Ko ngā tangata katoa, e manaakitia ana te whenua, o te Ao Whānui
(To all those who care for the lands of the wide world)
Ko ngā kaitiaki, e riterite ana, ngā whenua, huri rauna, i te Ao
(To the guardians of sustainability around the world)
Me tū tatou ki te werohia I ngā wero
(Stand and challenge the challenges)
I te ingoa o te tika o ngā mea katoa
(In the name of what is right with all things.)
Tēnā koutou katoa
(Greetings to you all)

Mr. President,

Mr. Secretary-General,

Friends,

I greet you in te reo Māori, language of the tangata whenua, or first people of Aotearoa New Zealand. I do so not just because it is the same way I would begin an address if I were at home, but because there are challenges we face as a world that I know no better way to express. Māori concepts like kaitiakitanga. The idea that each of us here today are guardians. Guardians of the land, of our environment and of our people. There is a simplicity to the notion of sovereign guardianship.

For decades we have assembled here under the assumption that we narrowly cooperate only on the issues that overtly impact on one another; issues like international trade rules, the law of the sea, or humanitarian access to war zones. The space in between has essentially been left to us. We, the political leaders of the world, have been the authors of our own domestic politics and policies. Decisions have been our own, and we have ultimately lived with the consequences.

But the world has changed.

Over time we have become increasingly interdependent. We see more and more often domestic decisions that have global ramifications.

Physical events have taught us that in obvious ways: oil spills that show no respect for maritime boundaries; nuclear accidents and testing, the impacts of which are never confined to the exact location in which they occur.

But our interdependence, our connection, runs so much deeper than that, and experiences in recent years should lead us all to question whether any of us ever truly operate in isolation anymore.

This is a question that we, the remote but connected nation of New Zealand, have been grappling with this year.

There are things that we are well known for in New Zealand. Green rolling hills, perfect you might say for hobbits to hide and for plenty of sheep to roam. We’re known for manaakitanga, or the pride we take in caring for our guests, so much so that it even extends to our most entrenched sporting rivals.

And now we are known for something else.

The 15th of March 2019.

The day an alleged terrorist undertook the most horrific attack on a place of worship, taking the lives of 51 innocent people, devastating our Muslim community and challenging our sense of who we are as a country.

Now there is no changing a nation’s history, but we can choose how it defines us. And in Aotearoa New Zealand, the people who lined up outside of mosques with flowers, the young people who gathered spontaneously in parks and open spaces in a show of solidarity, the thousands who stopped in silence to acknowledge the call to prayer seven days later, and the Muslim community who showed only love – these are the people who collectively decided that New Zealand would not be defined by an act of brutality and violence, but instead by compassion and empathy.

Make no mistake though, we do not claim to be a perfect nation. While we are a home of more than 200 ethnicities, that does not mean we are free from racism and discrimination. We have wounds from our own history that, 250 years on from the first encounters between Māori and Europeans, we continue to address.

But since the terrorist attack in New Zealand, we have had to ask ourselves many hard and many difficult questions.

One example sticks in my mind. It was only days after the shooting and I visited a mosque in our capital city. After spending some time with community leaders, I exited and walked across the car park where members of the Muslim community were gathered. And out of the corner of my eye I saw a young boy gesture to me. He was shy, almost retreating towards a barrier, but he also had something he clearly wanted to say. I quickly crouched down next to him.

He didn’t say his name or even hello, he simply whispered, “Will I be safe now?”

What does it take for a child to feel safe? As adults, we are quick to make the practical changes that will enable us to say that such a horrific act could never happen again. And we did that. Within 10 days of the attack we made a decision to change our gun laws and banned military style semi-automatic weapons and assault rifles in New Zealand. We have started on a second tranche of reforms to register weapons and to change our licensing regime.

These changes will help to make us safer. But when you’re a child, fear is not discrete, and it cannot be removed through legislative acts or decrees from Parliament.

Feeling safe means the absence of fear. Living free from racism, bullying and discrimination. Feeling loved, included and able to be exactly who you are. And to feel truly safe, those conditions need to be universal. No matter who you are, no matter where you come from, no matter where you live.

The young Muslim boy in Kilbirnie, New Zealand, wanted to know if I could grant him all of those things.

My fear is, that as a leader of a proudly independent nation, this is one thing that I cannot achieve alone. Not anymore.

In our borderless and technologically connected world, commentary on race, acts of discrimination based on religion, gender, sexuality or ethnicity – they are not neatly confined behind boundaries. They are felt globally.

The fact I received so many letters from Muslim children from around the world in the weeks after March 15 speaks to the power of connection. These children had no sense of distance. They may have never heard of New Zealand before March 15. They just saw an act of hatred against their community, and it felt close to them.

Whether it is acts of violence, language intended to incite fear of religious groups or assumptions about ethnicities to breed distrust and racism – these actions and utterances are as globalised as the movement of goods and services.

Children hear them.

Women hear them.

People of faith hear them.

Our rainbow communities hear them.

And so now, it’s our turn to stop and to listen. To accept that our words and actions have immeasurable consequences. And to speak not only like the whole world is listening, but with the responsibility of someone who knows a small child somewhere might be listening too.

The spaces in which we communicate, they are part of the challenge too, tough. In an increasingly online world we need to create spaces for the exchange of ideas, the sharing of technology and free speech, but also acknowledging the potential for this technology to be used to cause harm.

March 15 was a staggering example of such harm, and a deliberate effort to broadcast terror on a massive, viral scale over the internet. The alleged terrorist didn’t just take the lives of 51 people, he did it live on Facebook.

In the first 24 hours after the attack, Facebook took down 1.5 million copies of the livestream video. YouTube saw a copy of the video uploaded, at times, as fast as once every second during the same period.

The alleged terrorist used social media as a weapon. The attack demonstrated how the internet, a global commons with extraordinary power to do good, can be perverted and used as a tool for terrorists.

And so what happened in Christchurch, as well as being a profound tragedy, is also a complex and ongoing problem for the world. And it’s a problem we felt a sense of responsibility to do something about, so we sought to collaborate with the technology companies so integral to the solution.

Two months after the attacks, leaders gathered in Paris for the Christchurch Call, bringing together companies, countries and civil society, and committing to a range of actions to reduce the harm this content can cause. In doing so we have kept our focus on the deeper aim we all want: technology that unleashes human potential, not the worst in us.

Yesterday, I met with Call supporters to check in on our collective progress. We announced that a key tech industry institution will be reshaped to give effect to those commitments – and we launched a crisis response protocol to make sure that we can respond to such events in the future.

Neither New Zealand nor any other country could make these changes on their own. The tech companies couldn’t either. We are succeeding because we are working together, and for that unprecedented and powerful act of unity New Zealand says thank you.

The centrality of technology in our lives is not the only example though of our increasing interconnection and our reliance on one another if we are to respond to the challenges we face. There is perhaps no better example of our absolute interdependence than climate change.

When the United Nations Secretary-General visited the Pacific region this year, he saw first-hand how countries that have produced the fewest greenhouse gas emissions are now facing the most catastrophic effects. In his words: “To save the Pacific is to save the whole planet.”

In fact, seven out of the 15 most climate-affected nations in the world sit within the Pacific region. Places like Tuvalu, with a population of just over 11,000 people, barely contributes to global emissions but is paying the price for our collective inaction. Atolls so low lying that in weather events the water on either side of it can flow together and join at the narrowest points. Engulfed by the sea. Or Tokelau, a beautiful set of three atolls that can only be accessed by boat, where the children speak knowledgeably about climate change, knowing that unlike all of the challenges their self-reliant forbears have ever faced, this is the one that is completely and utterly in other people’s hands.

Now they have never met you, nor you them. But I can tell you that their expectations on all of us are high.

Meeting those expectations will require us to use every policy lever available – and, just like the Christchurch Call to Action, we need to work with partners inside and outside government to make change.

[…]

We are being asked to make decisions that are local, but with consequences that are global.

And yet, this is what climate change requires us to do. That is what historically our commitment to the United Nations Charter and the Universal Declaration of Human Rights asks us to do. It’s what standing up against acts of violence and discrimination asks us to do. Our globalised, borderless world asks us to be guardians not just for our people, but for all people.

There may have been a time when being unified under common challenges was an easier concept than it is today. But undeniably, we are living in a time where our greater reliance on one another has collided with a period of greater tribalism.

It would be wrong to assume that this is a new phenomenon. Research, in fact, has shown that humans are in fact so inclined to form natural tribes that if you put a completely unconnected diverse group of people into a room and flip a coin for each, those two groups will automatically form a suspicion of one another based on nothing more than heads vs tails.

Scientist and writer Robert Sapolsky recently reminded us that humans organise. Whether it’s class, race, country or coin flipping – there has always been a tendency to form us vs other.

But he also asks the question: what if we change what ‘us’ means? If instead of fierce nationalism or self-interest, we seek to form our tribes based on concepts that can and should be universal.

What if we no longer see ourselves based on what we look like, what religion we practise, or where we live. But by what we value. Humanity. Kindness. An innate sense of connection to each other.And a belief that we are guardians, not just of our home and our planet, but of each other.

We are borderless, but we are connected.

We are inherently different, but we have more that we share.

We may feel afraid, but as leaders we have the keys to create a sense of security, and a sense of hope.

We just need to choose.

Tātou tātou. No reira, tena koutou, tena kontou, tena koutou hatoa.
\end{document}
